\section{Conclusion}
Fact-equivalent attribute orders can determine whether a relevant product enters a neural retrieval pool, even when only that product changes. Two adapted visibility rewriters do not repair this problem: they amplify order sensitivity, often alter or omit source facts, and do not transfer consistently across datasets or encoders. Canonical preprocessing removes incoming-order variation, while a factuality guard makes the tested generation pipelines behave almost entirely like canonical text. Because canonicalization itself can reduce recall, robustness must be assessed jointly with relevant-item inclusion. The practical lesson is to treat product attributes as structured facts, verify every generated representation, and evaluate candidate loss before downstream ranking or agents are credited with improved visibility.
